\documentclass[openany, 10pt]{book}

%%%%%%%%%%%%%%%%%%%%%%%%%%%%%%%%%%%%%%%%%%%%%%%%%%%
%%% Find  10-readme-for-trm-repo.tex at:        %%%
%%% https://www.overleaf.com/read/bwdjvfjksvjy/ %%%
%%%%%%%%%%%%%%%%%%%%%%%%%%%%%%%%%%%%%%%%%%%%%%%%%%%

% >>> R?: marks a repo project setting number or important notice

% >>> R1: Configure chip-specific settings <<<
\usepackage{chip-spec-settings}

% Imports packages and global formatting settings
\usepackage{preamble-trm-repo}

% >>> R2: In file below, comment out the modules already included into
% this project to avoid warnings for multiple label definitions <<<
\myexternaldocument{temp-labels-en}

% Make text left-aligned and allow latex to break pages naturally, comment out for full justification
\raggedright
\raggedbottom

% Keep commented out!
% >>> Only uncomment to generate the list of added hypertargets <<<
%\toggletrue{addedhypertargets}


\begin{document}

\selectlanguage{english}
\setenglishformatting

% Initial checks make sure that all variables for cross-project shared
% content are configured correctly. If not, guidance will be provided.
%%%%%%%%%%%%%%%%%%%%%%%%%%%%%%%%%%%%%%%%%%%%%%%%%%%%%%%%%%%%%%%%%%%%%%%%%%
%%% For complete instructions:                                         %%%
%%% 1. Go to https://www.overleaf.com/project/60e2ed1aefb70e595778b04b %%%
%%% 2. Compile readme-general.tex                                      %%%
%%% 3. See Section "Initial Checks"                                    %%%
%%%%%%%%%%%%%%%%%%%%%%%%%%%%%%%%%%%%%%%%%%%%%%%%%%%%%%%%%%%%%%%%%%%%%%%%%%


%%% Warning for projects with undefined docprojtype %%%

% Check if the variable docprojtype has been added to the project
% to mark the project type for conditional inclusion of content, etc.
\ifcsname docprojtype\endcsname%

    % Check if the variable docprojtype is assigned a defined document project name
    \ifthenelse{%
      \equal{\docprojtype}{chip-datasheet} \OR
      \equal{\docprojtype}{trm} \OR
      \equal{\docprojtype}{hdg} \OR
      \equal{\docprojtype}{mod-datasheet} \OR
      \equal{\docprojtype}{errata}
    }{%
    % The variable docprojtype has been assigned a defined document project name
    % Proceed compilation.
    }{%
    % The variable docprojtype has NOT been assigned a defined document project name
    \begin{tiplisting}
      In your \texttt{preamble-....sty}, assign one of the following values to \texttt{docprojtype}, depending on your document type:
      \begin{itemize}
        \item \texttt{chip-datasheet}
        \item \texttt{trm}
        \item \texttt{hdg}
        \item \texttt{mod-datasheet}
        \item \texttt{errata}
      \end{itemize}
    \end{tiplisting}
    \errmessage{The variable 'docprojtype' has an unauthorized value}
    \stop
    }%

\else
    % The variable docprojtype has not been added to the project
    \begin{tiplisting}
      The code block below defines the variable \texttt{docprojtype} to mark document project type for conditional inclusion of content, etc.
      \begin{enumerate}
        \item \textbf{Insert the code block below} into your \texttt{preamble-....sty}.
        \item \textbf{Assign one of the following values} to \texttt{docprojtype}, depending on your document type:
	    \begin{itemize}
	        \item \texttt{chip-datasheet}
	        \item \texttt{trm}
	        \item \texttt{hdg}
	        \item \texttt{mod-datasheet}
	        \item \texttt{errata}
	    \end{itemize}
      \end{enumerate}

	    Example:

\begin{lstlisting}[numbers=none]
%%%%%%%%%%%%%%%%%%%%%%%%
%%% System Variables %%%
%%%%%%%%%%%%%%%%%%%%%%%%

% >>> Variable to mark document project type
%      for conditional inclusion of content, etc. <<<
\newcommand\docprojtype{<assign doc type>}
\end{lstlisting}

    \end{tiplisting}

    \errmessage{Define the variable 'docprojtype'}
    \stop

\fi% \docprojtype

\begin{titlepage}


\newgeometry{left=0cm,right=0cm,bottom=0cm}

\begin{flushleft}
{\color{myred}\rule{30pt}{2.7cm}}
    \hspace{0.2cm}
    \begin{minipage}[b]{18cm}
    {\fontsize{36}{48}\textbf{\chipname{}}}
    \vspace{\baselineskip}
    \vfill
    {\fontsize{28}{18}\textbf{\color{mygrayy}\GetTranslation{doctype}}}\smallskip
    \hspace{0.2cm}
    {\fontsize{18}{8}\textbf{\color{mygrayy}\GetTranslation{docversionprefix}\docversion{}}}
    \end{minipage}
\end{flushleft}

\vspace{48em}

\begin{flushright}
    \begin{tikzpicture}[remember picture,overlay]
        \node[anchor=north west, 
        xshift=16.5cm, 
        yshift=5.5cm]
        at (current page.south west) 
        {\includegraphics[height=2.2cm]{esp-logo-standard-vertical.pdf}};
    \end{tikzpicture}
    \begin{tikzpicture}[remember picture,overlay]
        \node[anchor=north west, 
            xshift=20cm, 
            yshift=5.5cm] 
            at (current page.south west) 
            {\color{myred}\rule{30pt}{2.2cm}};
    \end{tikzpicture}    
\end{flushright}

\vspace{9em}

\centerline{www.espressif.com}


\end{titlepage}

\restoregeometry
\clearpage

\clearpage
\thispagestyle{empty}

\setcounter{page}{2}

\subsection*{About This Document}

%\subsection*{Preliminary Release Note}

The {\bfseries \chipname{} \GetTranslation{doctype}} is targeted at developers working on low level software projects that use the \chipname{} SoC. It describes the hardware modules listed below for the \chipname{} SoC and other products in \chipname{} series. The modules detailed in this document provide an overview, list of features, hardware architecture details, any necessary programming procedures, as well as register descriptions.

\subsection*{Navigation in This Document}
Here are some tips on navigation through this extensive document:
\begin{itemize}
    \item \hyperref[tab:release-status]{Release Status at a Glance} on the very next page is a minimal list of all chapters from where you can directly jump to a specific chapter.
    \item Use the \textbf{Bookmarks} on the side bar to jump to any specific chapters or sections from anywhere in the document. Note this PDF document is configured to automatically display \textbf{Bookmarks} when open, which is necessary for an extensive document like this one. However, some PDF viewers or browsers ignore this setting, so if you don't see the \textbf{Bookmarks} by default, try one or more of the following methods:
    \begin{itemize}
        \item Install a PDF Reader Extension for your browser;
        \item Download this document, and view it with your local PDF viewer;
        \item Set your PDF viewer to always automatically display the \textbf{Bookmarks} on the left side bar when open.
    \end{itemize}
    \item Use the native \textbf{Navigation} function of your PDF viewer to navigate through the documents. Most PDF viewers support to go \textbf{Up}, \textbf{Down}, \textbf{Previous}, \textbf{Next}, \textbf{Back}, \textbf{Forward} and
    \textbf{Page} with buttons, menu, or hot keys.
    \item You can also use the built-in \textbf{GoBack} button on the upper right corner on each and every page to go back to the previous place before you click a link within the document. Note this feature may only work with some Acrobat-specific PDF viewers (for example, Acrobat Reader and Adobe DC) and browsers with built-in Acrobat-specific PDF viewers or extensions (for example, Firefox).
\end{itemize}

\clearpage
% % \thispagestyle{empty}

\setcounter{page}{3}

\subsection*{Release Status at a Glance}
\markboth{\it{Release Status at a Glance}}{\Navigation}

\begin{table}[H]
\centering
\label{tab:release-status}
\begin{tabular}{|c|L{5.8cm}|c|c|L{5.8cm}|c|}
\hline
\rowcolor{lightgray}
\textbf{No.} & \textbf{Chapter} & \textbf{Progress}  & \textbf{No.} & \textbf{Chapter} & \textbf{Progress} \\ \hline
\multicolumn{3}{|c|}{\textbf{Part \uppercase\expandafter{\romannumeral1}. Microprocessor and Master}} & 20    &  \nameref*{mod:rsa}    &    \hyperref[mod:rsa]{Published} \\\hline
1    &  \nameref*{mod:pie}    & \hyperref[mod:pie]{Published}   & 21    &   \nameref*{mod:hmac}    &    \hyperref[mod:hmac]{Published}\\\hline
2    &  \nameref*{mod:ulp}    &    \hyperref[mod:ulp]{Published}   & 22    &  \nameref*{mod:digsig}    &    \hyperref[mod:digsig]{Published}     \\\hline
3    &  \nameref*{mod:dma}    &    \hyperref[mod:dma]{Published}   & 23    &  \nameref*{mod:extmemencr}    &    \hyperref[mod:extmemencr]{Published}  \\\hline
\multicolumn{3}{|c|}{\textbf{Part \uppercase\expandafter{\romannumeral2}. Memory Organization}} & 24    &  \nameref*{mod:glitch}    &    \hyperref[mod:glitch]{Published} \\\hline
4    &  \nameref*{mod:sysmem}    &    \hyperref[mod:sysmem]{Published}  & 25    &    \nameref*{mod:rng}    &    \hyperref[mod:rng]{Published} \\\hline
5    &  \nameref*{mod:efuse}    &    \hyperref[mod:efuse]{Published}  & \multicolumn{3}{|c|}{\textbf{Part \uppercase\expandafter{\romannumeral5}. Connectivity Interface}}     \\\hline
\multicolumn{3}{|c|}{\textbf{Part \uppercase\expandafter{\romannumeral3}. System Component}} & 26    &  \nameref*{mod:uart}    &    \hyperref[mod:uart]{Published} \\\hline
6    &  \nameref*{mod:iomuxgpio}    &    \hyperref[mod:iomuxgpio]{Published}   & 27    &  \nameref*{mod:i2c}    &    \hyperref[mod:i2c]{Published}  \\\hline
7    &  \nameref*{mod:resclk}    &    \hyperref[mod:resclk]{Published}  & 28    &  \nameref*{mod:i2s}    &    \hyperref[mod:i2s]{Published}    \\\hline
8    &  \nameref*{mod:bootctrl}    &    \hyperref[mod:bootctrl]{Published} & 29    &   \nameref*{mod:camlcdctrl}    &    \hyperref[mod:camlcdctrl]{Published}  \\\hline
9    &  \nameref*{mod:intmtrx}    &    \hyperref[mod:intmtrx]{Published} & 30    &  \nameref*{mod:spi}    &    \hyperref[mod:spi]{Published}   \\\hline
10    &  \nameref*{mod:lowpowm}    &    \hyperref[mod:lowpowm]{Published}  & 31    &  \nameref*{mod:twai}    &    \hyperref[mod:twai]{Published}  \\\hline
11    &  \nameref*{mod:systimer}    &    \hyperref[mod:systimer]{Published}  & 32    &    \nameref*{mod:usbotg}    &    \hyperref[mod:usbotg]{Published} \\\hline
12    &  \nameref*{mod:timg}    &    \hyperref[mod:timg]{Published}  & 33    &  \nameref*{mod:usbserialjtag}    &    \hyperref[mod:usbserialjtag]{Published}    \\\hline
13    &  \nameref*{mod:wdt}    &    \hyperref[mod:wdt]{Published}   & 34    &  \nameref*{mod:sdhost}    &    \hyperref[mod:sdhost]{Published}  \\\hline
14    &  \nameref*{mod:xtwdt}    &    \hyperref[mod:xtwdt]{Published}   & 35    &  \nameref*{mod:ledpwm}    &    \hyperref[mod:ledpwm]{Published}   \\\hline
15    &  \nameref*{mod:permctrl}    &    \hyperref[mod:permctrl]{Published}  & 36    &  \nameref*{mod:mcpwm}    &    \hyperref[mod:mcpwm]{Published}  \\\hline
16    &  \nameref*{mod:wctrl}    &    \hyperref[mod:wctrl]{Published}  & 37    &  \nameref*{mod:rmt}    &    \hyperref[mod:rmt]{Published}   \\\hline
17    &  \nameref*{mod:sysreg}    &    \hyperref[mod:sysreg]{Published}  & 38   & \nameref*{mod:pcnt}    &    \hyperref[mod:pcnt]{Published}  \\\hline
\multicolumn{3}{|c|}{\textbf{Part \uppercase\expandafter{\romannumeral4}. Cryptography/Security Component}} & 39    &  \nameref*{mod:sensor}    &    \hyperref[mod:sensor]{Published}   \\\hline
18    &  \nameref*{mod:sha}    &    \hyperref[mod:sha]{Published}     & 38   & \nameref*{mod:pcnt}    &    \hyperref[mod:pcnt]{Published}  \\\hline
19    &  \nameref*{mod:aes}    &    \hyperref[mod:aes]{Published}  & \multicolumn{3}{|c|}{\textbf{Part \uppercase\expandafter{\romannumeral6}. Analog Signal Processing}}  \\\hline
20    &  \nameref*{mod:rsa}    &    \hyperref[mod:rsa]{Published}  & 39    &  \nameref*{mod:sensor}    &    \hyperref[mod:sensor]{Published}   \\\hline
\end{tabular}
\end{table}

\urlstyle{rm}

\begin{tiplisting}
    \vspace{-3.7em}
    \begin{longtable}[l]{ m{13.8cm}  p{2.5cm} }

        % table contents
        \makecell[l]{Check the link or the QR code to make sure that you use the latest version of this document:\\
            %\vspace{0.5em}\\
            %\href{\docpathlatestEN}{\docpathlatestEN}}
            \parbox{13.8cm}{\url{\docpathlatestEN}}}
        & \qrcode[height=2cm]{\docpathlatestEN}\\
    \end{longtable}
\end{tiplisting}

\setcounter{table}{0}



\tableofcontents
\clearpage
\listoftables
\clearpage
\listoffigures

% >>> R3: List all todo notes from the current project? <<<
% 'Yes' - uncomment the lines below; 'No' - comment out
% Don't forget to deactivate before publishing
%\todototoc
%\listoftodos
%{\let\clearpage\relax \listofdonetodos}


% Included modules
% Uncomment to include

%%%%%%%%%%%%%%%%%%%%%%%%%%%%
%%% Module Group: System %%%
%%%%%%%%%%%%%%%%%%%%%%%%%%%%

% \part{Microprocessor and Master}

% % Note: Updates may be necessary based on the specific SoC. For instance, Xtensa-based SoCs lack CPU descriptions, and certain SoCs may omit processor instruction extensions.

% This part covers the essential processing elements of the system, diving into the microprocessor architecture of the ultra-low-power processor. Details include CPU instruction extensions and controllers for Direct Memory Access (DMA).

% \subfileinclude{08-PIE__EN}
% \subfileinclude{26-ULP__EN}
% \subfileinclude{16-DMA__EN}

% \part{Memory Organization}

% This part provides insights into the system's memory structure, discussing the organization and mapping of RAM, ROM, eFuse, and external memories, offering a framework for understanding memory-related subsystems.

% \subfileinclude{09-SYSMEM__EN}
% \subfileinclude{07-eFuse__EN}

% \part{System Component}

% % Note: Tailor the description based on the specific components present in the SoC.

% Encompassing a range of system-level functionalities, this part describes components related to system boot, clocks, GPIO, timers, watchdogs, interrupt handling, low-power management, and various system registers.

% \subfileinclude{21-IOMUXGPIO__EN}
% \subfileinclude{02-RESCLK__EN}
% \subfileinclude{35-BOOTCTRL__EN}
% \subfileinclude{22-INTMTRX__EN}
% \subfileinclude{28-LOWPOWM__EN}
% \subfileinclude{34-SYSTIMER__EN}
% \subfileinclude{14-TIMG__EN}
% \subfileinclude{15-WDT__EN}
% \subfileinclude{37-XTWDT__EN}
% \subfileinclude{10-PERMCTRL__EN}
% \subfileinclude{38-WCTRL__EN}
% \subfileinclude{23-SYSREG__EN}

% \part{Cryptography/Security Component}

% Dedicated to security features, this part explores cryptographic accelerators like SHA and ECC. It also covers digital signatures, random number generation, and encryption/decryption algorithms, showcasing the SoC’s capabilities in cryptography and secure data processing.

% \subfileinclude{27-SHA__EN}
% \subfileinclude{30-AES__EN}
% \subfileinclude{31-RSA__EN}
% \subfileinclude{29-HMAC__EN}
% \subfileinclude{13-DIGSIG__EN}
% \subfileinclude{33-EXTMEMENCR__EN}
% \subfileinclude{36-GLITCH__EN}
% \subfileinclude{06-RNG__EN}


%%%%%%%%%%%%%%%%%%%%%%%%%%%%%%%%%
%%% Module Group: Peripherals %%%
%%%%%%%%%%%%%%%%%%%%%%%%%%%%%%%%%

\part{Connectivity Interface}

% Note: Update the description based on the actual components implemented in the SoC, e.g., some SoCs do not have RMT.

This part addresses the connectivity aspects of the system, describing components related to various communication interfaces like I2C, I2S, SPI, UART, USB, and more. The part also covers interfaces to generate signals used in remote control, motor control, LED control, etc.

% \subfileinclude{12-UART__EN}
\subfileinclude{01-I2C__EN}
\subfileinclude{03-I2S__EN}
% \subfileinclude{04-CAMLCDCTRL__EN}
% \subfileinclude{24-SPI__EN}
% \subfileinclude{32-TWAI__EN}
% \subfileinclude{11-USBOTG__EN}
% \subfileinclude{39-USBSERIALJTAG__EN}
% \subfileinclude{18-SDHOST__EN}
% \subfileinclude{19-LEDPWM__EN}
% \subfileinclude{17-MCPWM__EN}
% \subfileinclude{05-RMT__EN}
% \subfileinclude{20-PCNT__EN}

% \part{Analog Signal Processing}

% Note: Tailor the description based on the actual components implemented in the SoC.

% This part describes components related to analog-to-digital conversion, on-chip sensors, and features such as temperature sensing, demonstrating the system's capabilities in handling analog signals.

% \subfileinclude{25-SENSOR__EN}


%%%%%%%%%%%%%%%%
%%% Appendix %%%
%%%%%%%%%%%%%%%%

\part{Appendix}\label{appendix}

This part contains the following information starting from the next page:

\begin{itemize}
    \item \hyperref[sec:resources]{\textit{Related Documentation and Resources}}
    \item \hyperref[glossary]{\textit{Glossary}}
    \item \hyperref[programming-reserved-field]{\textit{Programming Reserved Register Field}}
    \item \hyperref[interrupt-config-registers]{\textit{Interrupt Configuration Registers}}
    \item \hyperref[revision-history]{\textit{Revision History}}
\end{itemize}


%%%%%%%%%%%%%%%%%%%%%%%%%%%%%%%%%%%%%%%%%%%
%%% Related Documentation and Resources %%%
%%%%%%%%%%%%%%%%%%%%%%%%%%%%%%%%%%%%%%%%%%%

\clearpage
%%%%%%%%%%%%%%%%%%%%%%%%%%%%%%%%%%%%%%%%%%%%%%%%%%%%%%%%%%%%%%%%%%%%%%%%%%
%%% For complete instructions:                                         %%%
%%% 1. Go to https://www.overleaf.com/project/60e2ed1aefb70e595778b04b %%%
%%% 2. Compile readme-general.tex                                      %%%
%%% 3. See Section "Add Section Related Documentation and Resources"   %%%
%%%%%%%%%%%%%%%%%%%%%%%%%%%%%%%%%%%%%%%%%%%%%%%%%%%%%%%%%%%%%%%%%%%%%%%%%%


% >>> Add chip name here <<<
% Always add '%' after the last chip name below (before the line '}')
\newcommand\definedchipnames{%
ESP32,
ESP32-S2,
ESP32-S3,
ESP32-C3,
ESP32-H2,
ESP8685,
ESP8684,
ESP32-C6,
ESP32-PICO,
ESP32-S3-PICO-1,
ESP8686,
ESP32-P4,
ESP32-C5,
ESP32-C61/ESP32-C31,
ESP-Shelly-C38F,
XMCC3ESP-FH%
}


% Code needed for the Introduced variable \definedchipnames
% to define all chip names in one place
% See https://tex.stackexchange.com/questions/642211/using-a-macro-to-pass-parameters-to-a-command/
% replace \docsvlist in \iftagged with \xdocsvlist
\patchcmd{\iftagged}{\docsvlist}{\xdocsvlist}{}{}
\patchcmd{\tagged}{\docsvlist}{\xdocsvlist}{}{}
\patchcmd{\untagged}{\docsvlist}{\xdocsvlist}{}{}
% \xdocsvlist fully expands its argument before passing it to \docsvlist
\newcommand{\xdocsvlist}[1]{%
  \expandafter\docsvlist\expandafter{\expanded{#1}}%
}


\phantomsection
\chapter*{Related Documentation and Resources}\label{sec:resources}
\markboth{\it{Related Documentation and Resources}}{\Navigation}
\addcontentsline{toc}{chapter}{Related Documentation and Resources} 


%%%%%%%%%%%%%%%%%%%%%%%%%%%%%%%%%%%%%%%%%%%%%%%%%
%%% Warning for projects with undefined chips %%%
%%%%%%%%%%%%%%%%%%%%%%%%%%%%%%%%%%%%%%%%%%%%%%%%%

% Show content below if undefined chip name encountered
\untagged{\definedchipnames}
{%
    \begin{tiplisting}
        Assign a defined chip name in your preamble > variable "chipseries"
    
	    Otherwise, define chip in 
	    \begin{itemize}
	        \item \texttt{resources-link-defs}
	        \item \texttt{resources-[doc-type]\_\_EN}
	        \item \texttt{resources-[doc-type]\_\_CN}
	    \end{itemize}
    \end{tiplisting}
}% \untagged{\definedchipnames}



%%%%%%%%%%%%%%%%%%%%%%%%%%%%%%%%%%%%%%%%%%%%%%%
%%% Content for projects with defined chips %%%
%%%%%%%%%%%%%%%%%%%%%%%%%%%%%%%%%%%%%%%%%%%%%%%

% Show content below if defined chip name encountered
\tagged{\definedchipnames}
{%

%%% Related Documentation %%%

\vspace{-0.5em}
\begin{minipage}{\textwidth}

\section*{Related Documentation}

\begin{itemize}
  % Include this link only if it is defined in resources-link-defs.tex
  \ifcsname linkchipdatasheetEN\endcsname%
    % Exclude this link if this section is added to a chip datasheet project
    \ifthenelse{\equal{\docprojtype}{chip-datasheet}}{}%
    {%
      \item \href{\linkprefix\linkchipdatasheetEN}{\chipseries{} Series Datasheet} --
        {\small Specifications of the \chipseries{} hardware.}%
    }%
  \fi%

  % Include this link only if it is defined in resources-link-defs.tex
  \ifcsname linkchiptrmEN\endcsname%
    % Exclude this link if this section is added to a TRM project
    \ifthenelse{\equal{\docprojtype}{trm}}{}%
    {%
      \item \href{\linkprefix\linkchiptrmEN}{\chipseries{} Technical Reference Manual} --
        {\small Detailed information on how to use the \chipseries{} memory and peripherals.}%
    }%
  \fi%

  % Include this link only if it is defined in resources-link-defs.tex
  \ifcsname linkchiphwdesignEN\endcsname%
    % Exclude this link if this section is added to a hardware design guidelines project
    \ifthenelse{\equal{\docprojtype}{hdg}}{}%
    {%
      \item \href{\linkprefix\linkchiphwdesignEN}{\chipseries{} Hardware Design Guidelines} --
        {\small Guidelines on how to integrate the \chipseries{} into your hardware product.}%
    }%
  \fi%

  % Include this link only if it is defined in resources-link-defs.tex
  \ifcsname linkchipecoEN\endcsname%
    \item \href{\linkprefix\linkchipecoEN}{\chipseries{} ECO and Workarounds for Bugs} --
        {\small Correction of \chipseries{} design errors.}
  \fi%
  
    % Include this link only if it is defined in resources-link-defs.tex
  \ifcsname linkchiperrataEN\endcsname%
   % Exclude this link if this section is added to an errata project
    \ifthenelse{\equal{\docprojtype}{errata}}{}%
    {%
    \item \href{\linkprefix\linkchiperrataEN}{\chipseries{} Series SoC Errata} --
        {\small Descriptions of known errors in \chipseries{} series of SoCs.}
        }%
  \fi%
  
  % If the variable \productcertificates is defined in a target project
  % append it to the link to filter available certificates
  \ifcsname productcertificates\endcsname%
    \item \textit{Certificates}\\
        \href{\linkprefix\linkcertificatesEN\productcertificates}{\linkprefix\linkcertificatesEN}
  \else
    \PackageWarning{Related Documentation and Resources}{Link extension for certificates not defined}
    \item \textit{Certificates}\\
        \href{\linkprefix\linkcertificatesEN}{\linkprefix\linkcertificatesEN}
  \fi%

  % Include this link only if it is defined in resources-link-defs.tex
  \ifcsname linkpcnEN\endcsname%
    \item \textit{\chipseries{} Product/Process Change Notifications (PCN)}\\
        \href{\linkprefix\linkpcnEN}{\linkprefix\linkpcnEN}
  \fi%

  % Include this link only if it is defined in resources-link-defs.tex
  \ifcsname linkadvisoryEN\endcsname%
    \item \textit{\chipseries{} Advisories} --
        {\small Information on security, bugs, compatibility, component reliability.}\\
        \href{\linkprefix\linkadvisoryEN}{\linkprefix\linkadvisoryEN}
  \fi%

    \item \textit{Documentation Updates and Update Notification Subscription}\\
        \href{\linkprefix\linkrecentlyupdatedEN}{\linkprefix\linkrecentlyupdatedEN}

\end{itemize}

\end{minipage}


%%% Developer Zone %%%

\begin{minipage}{\textwidth}

\section*{Developer Zone}

\begin{itemize}
  % Include this link only if it is defined in resources-link-defs.tex
  \ifcsname linkidfprogguideEN\endcsname%
    \item \href{\linkprefix\linkidfprogguideEN}{ESP-IDF Programming Guide for \chipseries{}} --
        {\small Extensive documentation for the ESP-IDF development framework.}
  \fi%

    \item \textit{ESP-IDF} and other development frameworks on GitHub.\\
        \href{\linkprefix\linkgithub}{\linkprefix\linkgithub}

    \item \textit{ESP32 BBS Forum} --
        {\small Engineer-to-Engineer (E2E) Community for Espressif products where you can post questions, share knowledge, explore ideas, and help solve problems with fellow engineers.}\\
        \href{\linkprefix\linkespforum}{\linkprefix\linkespforum}
        
    \item \textit{The ESP Journal} --
        {\small Best Practices, Articles, and Notes from Espressif folks.}\\
        \href{\linkprefix\linkespjournal}{\linkprefix\linkespjournal}

    \item See the tabs \textit{SDKs and Demos}, \textit{Apps}, \textit{Tools}, \textit{AT Firmware}.\\
        \href{\linkprefix\linksdkdemoEN}{\linkprefix\linksdkdemoEN}
\end{itemize}

\end{minipage}


%%% Products %%%

\begin{minipage}{\textwidth}

\section*{Products}

\begin{itemize}
    \item \textit{\chipseries{} Series SoCs} --
        {\small Browse through all \chipseries{} SoCs.}\\
        \href{\linkprefix\linkespsocsEN?id=\chipseries}{\linkprefix\linkespsocsEN?id=\chipseries}
    \item \textit{\chipseries{} Series Modules} --
        {\small Browse through all \chipseries{}-based modules.}\\
        \href{\linkprefix\linkespmodulesEN?id=\chipseries}{\linkprefix\linkespmodulesEN?id=\chipseries}
    \item \textit{\chipseries{} Series DevKits} --
        {\small Browse through all \chipseries{}-based devkits.}\\
        \href{\linkprefix\linkespdevkitsEN?id=\chipseries}{\linkprefix\linkespdevkitsEN?id=\chipseries}
    \item \textit{ESP Product Selector} --
        {\small Find an Espressif hardware product suitable for your needs by comparing or applying filters.}\\
        \href{\linkprefix\linkespproductselectorEN}{\linkprefix\linkespproductselectorEN}
\end{itemize}

\end{minipage}


%%% Contact Us %%%

\begin{minipage}{\textwidth}

\section*{Contact Us}

\begin{itemize}
    \item See the tabs \textit{Sales Questions}, \textit{Technical Enquiries}, \textit{Circuit Schematic \& PCB Design Review}, \textit{Get Samples} (Online stores), \textit{Become Our Supplier}, \textit{Comments \& Suggestions}.\\
        \href{\linkprefix\linksalesEN}{\linkprefix\linksalesEN}
\end{itemize}

\end{minipage}

}% \tagged{\definedchipnames}


\clearpage
\phantomsection
\chapter*{Glossary}\label{glossary}
\markboth{\it{Glossary}}{\Navigation}
\addcontentsline{toc}{chapter}{Glossary}

\phantomsection
\section*{Abbreviations for Peripherals}
\addcontentsline{toc}{section}{Abbreviations for Peripherals}

\begin{longtable}[c]{ R{3cm} p{12cm} }

AES     & AES (Advanced Encryption Standard) Accelerator \\
BOOTCTRL & Chip Boot Control \\ %newly added
DS      & Digital Signature\\ %newly added
DMA     & DMA (Direct Memory Access) Controller \\
eFuse   & eFuse Controller \\ 
% EMAC  & Ethernet MAC (delete it as it is not included in the name conventions of ESP32-S3-TRM)
HMAC    & HMAC (Hash-based Message Authentication Code) Accelerator \\ %newly added
I2C     & I2C (Inter-Integrated Circuit) Controller \\ %changed from  I²C 
I2S     & I2S (Inter-IC Sound) Controller \\ %changed from I²S
LEDC    & LED Control PWM (Pulse Width Modulation)\\ %changed from LED_PWM controller
MCPWM   & Motor Control PWM (Pulse Width Modulation) \\ 
%MMU    & Memory Management Unit (delete it as it is not included in the name conventions of ESP32-S3-TRM)
%MPU    & Memory Protection Unit (delete it as it is not included in the name conventions of ESP32-S3-TRM)
PCNT & Pulse Count Controller \\
% PERI  & Peripheral (Although it is not included in the name conventions, I am not sure whether it should be retained or not )
% PID   & Process Identifier (Although it is not included in the name conventions, I am not sure whether it should be retained or not  )
RMT     & Remote Control Peripheral \\
RNG     & Random Number Generator \\
RSA     & RSA (Rivest Shamir Adleman) Accelerator \\
% RTC   & Real Time Controller (Although it is not included in the name conventions, I am not sure whether it should be retained or not  )
SDHOST  & SD/MMC Host Controller \\ %changed from SDMMC to SDHOST
SHA     & SHA (Secure Hash Algorithm) Accelerator \\
SPI     & SPI (Serial Peripheral Interface) Controller \\
SYSTIMER & System Timer \\ %newly added
TIMG    & Timer Group \\
TWAI    & Two-wire Automotive Interface \\ %newly added
UART    & UART (Universal Asynchronous Receiver-Transmitter) Controller \\ 
ULP Coprocessor & Ultra-low-power Coprocessor \\ %changed from ULP Co-processor to ULP Coprocessor
USB OTG & USB On-The-Go \\ %newly added
WDT & Watchdog Timers \\ 
% XTWDT & XTAL32 Watchdog Timer\\ remove this line 

\end{longtable}

%%% Registers %%%

% Include sections
% - Abbreviations Related to Registers
% - Access Types for Registers

%%% Abbreviations Related to Registers %%%

\phantomsection
\section*{Abbreviations Related to Registers}\label{glossary-abbr-reg}
\addcontentsline{toc}{section}{Abbreviations Related to Registers}

\begin{longtable}[c]{ R{3cm} p{12cm} }

REG     & \textbf{Register}. \\
SYSREG  & \textbf{System registers} are a group of registers that control system reset, memory, clocks, software interrupts, power management, clock gating, etc. \\
ISO     & \textbf{Isolation}. If a peripheral or other chip component is powered down, the pins, if any, to which its output signals are routed will go into a floating state. ISO registers isolate such pins and keep them at a certain determined value, so that the other non-powered-down peripherals/devices attached to these pins are not affected. \\
NMI     & \textbf{Non-maskable interrupt} is a hardware interrupt that cannot be disabled or ignored by the CPU instructions. Such interrupts exist to signal the occurrence of a critical error. \\
W1TS    & Abbreviation added to names of registers/fields to indicate that such register/field should be used to set a field in a corresponding register with a similar name. For example, the register \texttt{GPIO\_ENABLE\_W1TS\_REG} should be used to set the corresponding fields in the register \texttt{GPIO\_ENABLE\_REG}. \\
W1TC    & Same as \textit{W1TS}, but used to clear a field in a corresponding register. \\
% FPU   &
% FPD   &
% PLLA  &
% BBPLL &

\end{longtable}


%%% Access Types for Registers %%%

% \clearpage
\phantomsection
\section*{Access Types for Registers}\label{glossary-access-types}
\addcontentsline{toc}{section}{Access Types for Registers}

Sections \textit{Register Summary} and \textit{Register Description} in TRM chapters specify access types for registers and their fields.

Most frequently used access types and their combinations are as follows:

\begin{multicols}{4}
  \begin{itemize}
    \item RO
    \item WO
    \item WT
    \item R/W
    \item WL
    \item R/W/SC
    \item R/W/SS
    \item R/W/SS/SC
    \item R/WC/SS
    \item R/WC/SC
    \item R/WC/SS/SC
    \item R/WS/SC
    \item R/WS/SS
    \item R/WS/SS/SC
    \item R/SS/WTC
    \item R/SC/WTC
    \item R/SS/SC/WTC
    \item RF/WF
    \item R/SS/RC
    \item varies
  \end{itemize}
\end{multicols}

Descriptions of all access types are provided below.


\begin{longtable}[c]{ R{1cm} p{13cm} }

% table contents
R & \textbf{Read.} 
	User application can read from this register/field; usually combined with other access types. \\
RO & \textbf{Read only.} 
	User application can only read from this register/field. \\
HRO & \textbf{Hardware Read Only.} 
	Only hardware can read from this register/field; used for storing default settings for variable parameters. \\
W & \textbf{Write.} 
	User application can write to this register/field; usually combined with other access types. \\
WO & \textbf{Write only.} 
	User application can only write to this register/field. \\

SS & \textbf{Self set.} 
	On a specified event, hardware automatically writes 1 to this register/field; used with 1-bit fields. \\
SC & \textbf{Self clear.} 
	On a specified event, hardware automatically writes 0 to this register/field; used with 1-bit and multi-bit fields. \\
SM & \textbf{Self modify.} 
	On a specified event, hardware automatically writes a specified value to this register/field; used with multi-bit fields. \\
SU & \textbf{Self update.} 
	On a specified event, hardware automatically updates this register/field; used with multi-bit fields. \\

RS & \textbf{Read to set.} 
	If user application reads from this register/field, hardware automatically writes 1 to it. \\
RC & \textbf{Read to clear.} 
	If user application reads from this register/field, hardware automatically writes 0 to it. \\
RF & \textbf{Read from FIFO.} 
	If user application writes new data to FIFO, the register/field automatically reads it. \\

WF & \textbf{Write to FIFO.} 
	If user application writes new data to this register/field, it automatically passes the data to FIFO via APB bus. \\

WS & \textbf{Write any value to set.} 
	If user application writes to this register/field, hardware automatically sets this register/field. \\
W1S & \textbf{Write 1 to set.} 
	If user application writes 1 to this register/field, hardware automatically sets this register/field. \\
W0S & \textbf{Write 0 to set.} 
	If user application writes 0 to this register/field, hardware automatically sets this register/field. \\

WC & \textbf{Write any value to clear.} 
	If user application writes to this register/field, hardware automatically clears this register/field. \\
W1C & \textbf{Write 1 to clear.} 
	If user application writes 1 to this register/field, hardware automatically clears this register/field. \\
W0C & \textbf{Write 0 to clear.} 
	If user application writes 0 to this register/field, hardware automatically clears this register/field. \\

WT & \textbf{Write 1 to trigger an event.} 
	If user application writes 1 to this field, this action triggers an event (pulse in the APB bus) or clears a corresponding WTC field (see WTC). \\
WTC & \textbf{Write to clear.} 
	Hardware automatically clears this field if user application writes 1 to the corresponding WT field (see WT). \\

W1T & \textbf{Write 1 to toggle.} 
	If user application writes 1 to this field, hardware automatically inverts the corresponding field; otherwise - no effect. \\
W0T & \textbf{Write 0 to toggle.} 
	If user application writes 0 to this field, hardware automatically inverts the corresponding field; otherwise - no effect. \\

WL & \textbf{Write if a lock is deactivated.} 
	If the lock is deactivated, user application can write to this register/field. \\

varies & \textbf{The access type varies.} 
	 Different fields of this register might have different access types. \\

\end{longtable}


% \subsection*{Abbreviations for Registers}
% \phantomsection
% \addcontentsline{toc}{subsection}{Abbreviations for Registers}

% \begin{longtable}[c]{ R{3cm} p{12cm} }

% ISO     & Isolation. When a module is power down, its output pins will be stuck in unknown state (some middle voltage). "ISO" registers will control to isolate its output pins to be a determined value, so it will not affect the status of other working modules which are not power down.\\
% NMI     & Non-maskable interrupt. \\
% REG     & Register.\\
% R/W     & Read/write. Software can read and write to these bits. \\
% RO      & Read-only. Software can only read these bits. \\
% SYSREG  & System Registers \\ %newly added
% WO      & Write-only. Software can only write to these bits.\\

% \end{longtable}


\clearpage
\phantomsection
\chapter*{Programming Reserved Register Field}\label{programming-reserved-field}
\markboth{\it{Programming Reserved Register Field}}{\Navigation}
\addcontentsline{toc}{chapter}{Programming Reserved Register Field}

\phantomsection
\section*{Introduction}
\addcontentsline{toc}{section}{Introduction}

A field in a register is reserved if the field is not open to users, or produces unpredictable results if configured to values other than defaults.

\phantomsection
\section*{Programming Reserved Register Field}
\addcontentsline{toc}{section}{Programming Reserved Register Field}

The reserved fields should not be modified. It is not possible to write only part of a register since registers must always be written as a whole. As a result, to write an entire register that contains reserved fields, you can choose one of the following two options:

\begin{enumerate}
    \item Read the value of the register, modify only the fields you want to configure and then write back the value so that reserved fields are untouched.

\noindent OR

\item Modify only the fields you want to configure and write back the default value of the reserved fields. The default value of a field is provided in the "Reset" line of a register diagram. For example, the default value of Field\_A in \hyperref[fig:registerX]{Register X} is 1.
\end{enumerate}

\begin{register}{H}{Register X}{Address}\label{fig:registerX}
\regfield{(reserved)}{12}{20}{000000000000}%
\regfield{Field\_C}{4}{16}{{0000}}%
\regfield{(reserved)}{14}{2}{00000000000000}%
\regfield{Field\_B}{1}{1}{{0}}%
\regfield{Field\_A}{1}{0}{{1}}%
\reglabel{Reset}\regnewline%

\end{register}

Suppose you want to set Field\_A, Field\_B, and Field\_C of \hyperref[fig:registerX]{Register X} to 0x0, 0x1, and 0x2, you can:

\begin{itemize}
    \item Use option 1 and fill in the reserved fields with the value you have just read. Suppose the register reads as 0x0000\_0003. Then, you can modify the fields you want to configure, thus writing 0x0002\_0002 to the register.
    \item Use option 2 and fill in the reserved fields with their defaults, thus writing 0x0002\_0002 to the register.
\end{itemize}


\clearpage
\phantomsection
\chapter*{Interrupt Configuration Registers}\label{interrupt-config-registers}
\markboth{\it{Interrupt Configuration Registers}}{\Navigation}
\addcontentsline{toc}{chapter}{Interrupt Configuration Registers}

Generally, the peripherals' internal interrupt sources can be configured by the following common set of registers:

\begin{itemize}
    \item \textbf{RAW} (Raw Interrupt Status) register: This register indicates the raw interrupt status. Each bit in the register represents a specific internal interrupt source. When an interrupt source triggers, its RAW bit is set to 1.

    \item \textbf{ENA} (Enable) register: This register is used to enable or disable the internal interrupt sources. Each bit in the ENA register corresponds to an internal interrupt source.

    By manipulating the ENA register, you can mask or unmask individual internal interrupt source as needed. When an internal interrupt source is masked (disabled), it will not generate an interrupt signal, but its value can still be read from the RAW register.

    \item \textbf{ST} (Status) register: This register reflects the status of enabled interrupt sources. Each bit in the ST register corresponds to a specific internal interrupt source. The ST bit being 1 means that both the corresponding RAW bit and ENA bit are 1, indicating that the interrupt source is triggered and not masked. The other combinations of the RAW bit and ENA bit will result in the ST bit being 0.

    The configuration of ENA/RAW/ST registers is shown in Table \ref{tab:int-raw-st-ena}.

    \item \textbf{CLR} (Clear) register: The CLR register is responsible for clearing the internal interrupt sources. Writing 1 to the corresponding bit in the CLR register clears the interrupt source.
\end{itemize}

\begin{longtable}{ | l | l | l | }
\caption{Configuration of ENA/RAW/ST Registers} \label{tab:int-raw-st-ena}
\\ \hline
\rowcolor{lightgray}
\textbf{ENA Bit Value} & \textbf{RAW Bit Value}  &  \textbf{ST Bit Value}\\ \hline
0 & Ignored & 0 \\ \hline
\multirow{2}*{1} &  0 & 0 \\ \cline{2-3}
 &  1 & 1 \\ \hline
\end{longtable}

\clearpage
\phantomsection
\chapter*{Revision History}\label{revision-history}
\markboth{\it{Revision History}}{\Navigation}
\addcontentsline{toc}{chapter}{Revision History}

\begin{longtable}[l]{ | m{2cm} | m{1.5cm} | m{12cm} | } \hline

    \rowcolor{lightgray}
    Date & Version & \hypertarget{Release notes}{Release notes}\\\hline
    \endfirsthead

    \multicolumn{3}{c}{\bfseries Cont'd from previous page}\\\hline
    \rowcolor{lightgray}
    Date & Version & Release notes\\\hline
    \endhead

    \multicolumn{3}{r}{\textbf{Cont'd on next page}}\\
    \endfoot

    \endlastfoot

    % table contents
    % 2022-xx-xx     & v1.1  &
    % Updated the following chapters:
    % \begin{itemize}
    %     \item Chapter \ref{mod:lowpowm} \textit{\nameref{mod:lowpowm}}
    %     \item Chapter \ref{mod:usbotg} \textit{\nameref{mod:usbotg}}
    % \end{itemize}
    % \\\hline
    2024- & v1.6 &
    Updated the following chapters:
        \begin{itemize}
            \item Chapter \ref{mod:lowpowm} \textit{\nameref{mod:lowpowm}}:
                \begin{itemize}
                    \item Updated the description of predefined power modes %TRMC6-42193
                    \item Added a note about EXT1 under Table \ref{tab:lpm-wakeup-source} %DOC-8414
                \end{itemize}
            \item Chapter \ref{mod:iomuxgpio} \textit{\nameref{mod:iomuxgpio}}: Updated the description of Hold feature %DOC-7326
            \item Chapter \ref{mod:sensor} \textit{\nameref{mod:sensor}}: Removed descriptions about the internal voltage/signal in SAR ADC2 measurement %DOC-7966
        \end{itemize}
    \\\hline
    2024-4-18 & v1.5 &
    Updated the following chapters:
        \begin{itemize}
            \item Chapter \ref{mod:dma} \textit{\nameref{mod:dma}}: Updated the descriptions of suc\_eof and the EOF flag %DOC-5478
            \item Chapter \ref{mod:iomuxgpio} \textit{\nameref{mod:iomuxgpio}}: Updated the drive strength of GPIO17 and GPIO18 %DOC-7383
            \item Chapter \ref{mod:digsig} \textit{\nameref{mod:digsig}} and Chapter \ref{mod:bootctrl} \textit{\nameref{mod:bootctrl}}: Fixed some typos %DOC-7472 %DOC-7629
            \item Chapter \ref{mod:uart} \textit{\nameref{mod:uart}}: Updated descriptions about clearing the wake\_up signal %DOC-7076
            \item Chapter \ref{mod:ledpwm} \textit{\nameref{mod:ledpwm}}: Updated the lowest resolution in Table \ref{tab:ledc-common-freq-res}  %DOC-7304
        \end{itemize}
    \\\hline
    2024-01-30      & v1.4  &
    Updated the following chapters:
    \begin{itemize}
        \item Chapter \ref{mod:pie} \textit{\nameref{mod:pie}}: Added data exchange instruction list to Section \ref{chp:pie_instruction_simple_introduction}; Added descriptions for LD.QR, ST.QR, and MV.QR instructions; Fixed two typos; %DOC-6202;%DOC-5679;%DOC-6560
        \item Chapter \ref{mod:efuse} \textit{\nameref{mod:efuse}}: Updated the description of \hyperref[fielddesc:EFUSEPINPOWERSELECTION]{EFUSE\_PIN\_POWER\_SELECTION} %DOC-5591
        \item Chapter \ref{mod:iomuxgpio} \textit{\nameref{mod:iomuxgpio}}: Removed Debug Assist from Table \ref{tab:sysmem-base-address}; Updated the description in Section \ref{PadHold} and Table \ref{tab:io_mux_reg_summary} %DOC-5892 %DOC-6733 %DOC-6563
        \item Chapter \ref{mod:bootctrl} \textit{\nameref{mod:bootctrl}}: Added SPI Download Boot mode and renamed Download Boot mode to Joint Download mode in Section \ref{sec:bootModeControl}; Provided more details about how FUSE\_DIS\_FORCE\_DOWNLOAD and EFUSE\_DIS\_DOWNLOAD\_MODE control chip boot mode %DOC-5916 %DOC-5892
        \item Chapter \ref{mod:intmtrx} \textit{\nameref{mod:intmtrx}}: Removed the ASSIST\_DEBUG\_INTR interrupt source %DOC-6562
        \item Chapter \ref{mod:lowpowm} \textit{\nameref{mod:lowpowm}}: Updated the description of register \hyperref[fielddesc:RTCCNTLWDTWKEY]{RTC\_CNTL\_WDT\_WKEY}%DOC-6760
        \item Chapter \ref{mod:timg} \textit{\nameref{mod:timg}}: Updated the description of \hyperref[fielddesc:TIMGWDTCLKPRESCALE]{TIMG\_WDT\_CLK\_PRESCALE} % DOC-5697
        \item Chapter \ref{mod:permctrl} \textit{\nameref{mod:permctrl}}: removed the description of access configuration for Debug Assist  %DOC-6562
        \item Chapter \ref{mod:sysreg} \textit{\nameref{mod:sysreg}}: Improved the description of \hyperref[fielddesc:SYSTEMCONTROLCORE1MESSAGE]{SYSTEM\_CONTROL\_CORE\_1\_MESSAGE}% DOC-5769
        \item Chapter \ref{mod:uart} \textit{\nameref{mod:uart}}: Updated the number of rising edges required to generate the wake\_up signal %DOC-7076
        \item Chapter \ref{mod:i2c} \textit{\nameref{mod:i2c}}: Updated the descriptions of \hyperref[regdesc:I2CCOMD0REG]{I2C\_COMD0\_REG}, \hyperref[fielddesc:I2CSDAFORCEOUT]{I2C\_SDA\_FORCE\_OUT} and \hyperref[fielddesc:I2CSCLFORCEOUT]{I2C\_SCL\_FORCE\_OUT}  %DOC-6428, DOC-6511
        \item Chapter \ref{mod:mcpwm} \textit{\nameref{mod:mcpwm}}: Added one note about Count-Up-Down mode configuration %DOC-7080
    \end{itemize}
    Added Section \hyperref[interrupt-config-registers]{\textit{Interrupt Configuration Registers}}

    \\\hline
    2023-07-04 & v1.3 &
    Updated the following chapters:
    \begin{itemize}
        \item Chapter \ref{mod:ulp} \textit{\nameref{mod:ulp}}: Added op\_code for RTC I2C instructions; Added more register fields into \hyperref[regdesc:SENSSARI2CCTRLREG]{SENS\_SAR\_I2C\_CTRL\_REG}; Fixed a typo; Added Section Interrupt Handling %DOC-4933 DOC-5089
        \item Chapter \ref{mod:dma} \textit{\nameref{mod:dma}}: Updated the descriptions of the GDMA\_IN\_SUC\_EOF\_CH\regindex{n}\_INT interrupt and the GDMA\_INLINK\_DSCR\_ADDR\_CH\regindex{n} field % DOC-4730
        \item Chapter \ref{mod:sysmem} \textit{\nameref{mod:sysmem}}: Updated Table \ref{tab:sysmem-base-address} and Figure \ref{fig:sysmem-address-mapping_structure} %DOC-5109
        \item Chapter \ref{mod:iomuxgpio} \textit{\nameref{mod:iomuxgpio}}: Updated the description of  \hyperref[fielddesc:IOMUXPADPOWERCTRL]{IO\_MUX\_PAD\_POWER\_CTRL} %DOC-5593
        \item Chapters \ref{mod:permctrl} \textit{\nameref{mod:permctrl}} and \ref{mod:sysreg} \textit{\nameref{mod:sysreg}}: Updated register prefix APB\_CTRL to SYSCON % DOC-5455
        \item Chapter \ref{mod:uart} \textit{\nameref{mod:uart}}: Updated Figure \nameref{fig:uart-arch-overview}, Figure \nameref{fig:one-uart-structure}, and Figure \nameref{fig:uart-flow-control}; Updated the maximum length of stop bits and related descriptions %DOC-4826
        \item Chapter \ref{mod:usbotg} \textit{\nameref{mod:usbotg}}: Added a note on how to access USB OTG registers %DOC-5008
        \item Chapter \ref{mod:usbserialjtag} \textit{\nameref{mod:usbserialjtag}}: Added the specific pull-up values configured by the \hyperref[fielddesc:USBSERIALJTAGPULLUPVALUE]{USB\_SERIAL\_JTAG\_PULLUP\_VALUE} bit
        \item Chapter \ref{mod:sensor} \textit{\nameref{mod:sensor}}: Updated the pattern table indexes in Figure \ref{fig:pattern-tab-1} and Figure \ref{fig:pattern-tab-2}
    \end{itemize}
    Added Section \hyperref[programming-reserved-field]{\textit{Programming Reserved Register Field}}
    \\\hline
    2023-03-02        & v1.2  &
    Updated the following chapters:
    \begin{itemize}
        \item Chapter \ref{mod:sysmem} \textit{\nameref{mod:sysmem}}: Updated the description for Internal ROM 1 in Section \ref{other:sysmem-internal-ROM-1} \textit{\nameref{other:sysmem-internal-ROM-1}}
        %\item Chapter \ref{mod:efuse} \textit{\nameref{mod:efuse}}: Fixed a typo %DOC-4232 (Mo: included in "other minor updates")
        \item Chapter \ref{mod:lowpowm} \textit{\nameref{mod:lowpowm}}: Removed UART as a reject to sleep cause %DOC-4335  %DOC-3662
        \item Chapter \ref{mod:timg} \textit{\nameref{mod:timg}}: Updated the procedures to read the timer's value  %DOC-3634
        \item Chapter \ref{mod:sysreg} \textit{\nameref{mod:sysreg}}: Added descriptions about the \hyperref[regdesc:SYSTEMPERIPRSTEN0REG]{SYSTEM\_PERIP\_RST\_EN0\_REG} and \hyperref[regdesc:SYSTEMPERIPCLKEN0REG]{SYSTEM\_PERIP\_CLK\_EN0\_REG} registers%DOC-4548
        \item Chapter \ref{mod:uart} \textit{\nameref{mod:uart}}: Added descriptions about the break condition %DOC-3633
        \item Chapter \ref{mod:usbserialjtag} \textit{\nameref{mod:usbserialjtag}}: Added descriptions about IO pad state change in USB-OTG download mode%doc-4373
        \item Chapter \ref{mod:ledpwm} \textit{\nameref{mod:ledpwm}}: Added the formula to calculate duty cycle resolution and Table  \textit{\nameref{tab:ledc-common-freq-res}} %DOC-3480
        \item Chapter \ref{mod:sensor} \textit{\nameref{mod:sensor}}: Added descriptions about the \hyperref[regdesc:APBSARADCAPBSARADC1DATASTATUSREG]{APB\_SARADC\_APB\_SARADC1\_DATA\_STATUS\_REG} register, updated the description about the \hyperref[fielddesc:SENSFORCEXPDSAR]{SENS\_FORCE\_XPD\_SAR} field, and added a note about the touch sensor limitation
    \end{itemize}
    Other minor updates
    \\\hline
    2022-10-31         & v1.1  &
    Updated the following chapters:
    \begin{itemize}
        \item Chapter \ref{mod:pie} \textit{\nameref{mod:pie}}
        \item Chapter \ref{mod:efuse} \textit{\nameref{mod:efuse}}
        \item Chapter \ref{mod:sensor} \textit{\nameref{mod:sensor}}%DOC-3564
    \end{itemize}
    Updated the \hyperref[glossary]{Glossary} section
    \\\hline
    2022-09-08         & v1.0  &
    Added the following chapter:
    \begin{itemize}
        \item Chapter \ref{mod:pie} \textit{\nameref{mod:pie}}
    \end{itemize}
    Updated the following chapters:
    \begin{itemize}
        \item Chapter \ref{mod:intmtrx} \textit{\nameref{mod:intmtrx}}
        \item Chapter \ref{mod:efuse} \textit{\nameref{mod:efuse}}
        \item Chapter \ref{mod:lowpowm} \textit{\nameref{mod:lowpowm}} %doc3469
        \item Chapter \ref{mod:rng} \textit{\nameref{mod:rng}} % doc-3469
        \item Chapter \ref{mod:usbotg} \textit{\nameref{mod:usbotg}} %
    \end{itemize}
    \\\hline
    2022-08-04     & v0.8  &
    Added the following chapters:
    \begin{itemize}
        \item Chapter \ref{mod:i2s} \textit{\nameref{mod:i2s}}
        \item Chapter \ref{mod:spi} \textit{\nameref{mod:spi}}
    \end{itemize}
    Updated the following chapters:
    \begin{itemize}
        \item Chapter \ref{mod:efuse} \textit{\nameref{mod:efuse}} %doc-3121
        \item Chapter \ref{mod:lowpowm} \textit{\nameref{mod:lowpowm}} %register description about w1ts/w1tc
    \end{itemize}
    \\\hline
    2022-06-30     & v0.7  &
    Added the following chapters:
    \begin{itemize}
        \item Chapter \ref{mod:sensor} \textit{\nameref{mod:sensor}}
        \item Chapter \ref{mod:lowpowm} \textit{\nameref{mod:lowpowm}}
    \end{itemize}
    Updated the following chapters:
    \begin{itemize}
        \item Chapter \ref{mod:sysmem} \textit{\nameref{mod:sysmem}}
        \item Updated clock names:
        \begin{itemize}
        \item FOSC\_CLK: renamed as RC\_FAST\_CLK
        \item FOSC\_DIV\_CLK: renamed as RC\_FAST\_DIV\_CLK
        \item RTC\_CLK: renamed as RC\_SLOW\_CLK
        \item SLOW\_CLK: renamed as RTC\_SLOW\_CLK
        \item FAST\_CLK: renamed as RTC\_FAST\_CLK
        \item PLL\_80M\_CLK: renamed as PLL\_F80M\_CLK
        \item PLL\_160M\_CLK: renamed as PLL\_F160M\_CLK
        \item PLL\_240M\_CLK: renamed as PLL\_D2\_CLK%doc-3141
    \end{itemize}
    \end{itemize}
    \\\hline
    2022-06-01     & v0.6  &
    Updated the following chapters:
    \begin{itemize}
        \item Chapter \ref{mod:dma} \textit{\nameref{mod:dma}}
        \item Chapter \ref{mod:efuse} \textit{\nameref{mod:efuse}} %DOC-2901
        \item Chapter \ref{mod:bootctrl} \textit{\nameref{mod:bootctrl}}
    \end{itemize}
    \\\hline
    2022-04-06     & v0.5  &
    Added the following chapters:
    \begin{itemize}
        \item Chapter \ref{mod:permctrl} \textit{\nameref{mod:permctrl}}
        \item Chapter \ref{mod:wctrl} \textit{\nameref{mod:wctrl}}
    \end{itemize}
    Updated the following chapter:
    \begin{itemize}
        \item Chapter \ref{mod:sysmem} \textit{\nameref{mod:sysmem}}
    \end{itemize}
    \\\hline

    2022-02-24    & v0.4  &
    Added Chapter \ref{mod:camlcdctrl} \textit{\nameref{mod:camlcdctrl}}

    Updated the following chapters:
        \begin{itemize}
            \item Chapter \ref{mod:ulp} \textit{\nameref{mod:ulp}} %DOC-2517, DOC-2516, DOC-2710, DOC-2611
            \item Chapter \ref{mod:efuse} \textit{\nameref{mod:efuse}}
            \item Chapter \ref{mod:iomuxgpio} \textit{\nameref{mod:iomuxgpio}} %DOC-2730
            \item Chapter \ref{mod:resclk} \textit{\nameref{mod:resclk}}
            \item Chapter \ref{mod:bootctrl} \textit{\nameref{mod:bootctrl}} %DOC-2569
            \item Chapter \ref{mod:usbotg} \textit{\nameref{mod:usbotg}}
             \end{itemize}
    \\\hline
    2021-12-16 & v0.3 &
    Added the following chapters:
        \begin{itemize}
            \item Chapter \ref{mod:ulp} \textit{\nameref{mod:ulp}}
            \item Chapter \ref{mod:dma} \textit{\nameref{mod:dma}}
            \item Chapter \ref{mod:systimer} \textit{\nameref{mod:systimer}}
            \item Chapter \ref{mod:glitch} \textit{\nameref{mod:glitch}}
            \item Chapter \ref{mod:i2c} \textit{\nameref{mod:i2c}}
            \item Chapter \ref{mod:mcpwm} \textit{\nameref{mod:mcpwm}}
            \item Chapter \ref{mod:rmt} \textit{\nameref{mod:rmt}}
        \end{itemize}
    Updated the following chapters:
        \begin{itemize}
            \item Chapter \ref{mod:efuse} \textit{\nameref{mod:efuse}}
            \item Chapter \ref{mod:rsa} \textit{\nameref{mod:rsa}}
            \item Chapter \ref{mod:hmac} \textit{\nameref{mod:hmac}}
            \item Chapter \ref{mod:digsig} \textit{\nameref{mod:digsig}}
        \end{itemize}
    \\\hline

    2021-09-30      & v0.2 &

            Added the following chapters:
            \begin{itemize}
                \item Chapter \ref{mod:sysreg} \textit{\nameref{mod:sysreg}}
                \item Chapter \ref{mod:hmac} \textit{\nameref{mod:hmac}}
                \item Chapter \ref{mod:extmemencr} \textit{\nameref{mod:extmemencr}}
                \item Chapter \ref{mod:uart} \textit{\nameref{mod:uart}}
                \item Chapter \ref{mod:usbserialjtag} \textit{\nameref{mod:usbserialjtag}}
            \end{itemize}
            Updated the following Chapters:
            \begin{itemize}
                \item Chapter \ref{mod:efuse} \textit{\nameref{mod:efuse}}
                \item Chapter \ref{mod:twai} \textit{\nameref{mod:twai}}
            \end{itemize}

            \\\hline
    2021-07-09 & v0.1 & Preliminary release\\\hline
\end{longtable}


\clearpage
%%%%%%%%%%%%%%%%%%%%%%%%%%%%%%%%%%%%%%%%%%%%%%%%%%%%%%%%%%%%%%%%%%%%%%%%%%
%%% For complete instructions:                                         %%%
%%% 1. Go to https://www.overleaf.com/project/60e2ed1aefb70e595778b04b %%%
%%% 2. Compile readme-general.tex                                      %%%
%%% 3. See Section "Updating Disclaimer and Copyright Notice"          %%%
%%%%%%%%%%%%%%%%%%%%%%%%%%%%%%%%%%%%%%%%%%%%%%%%%%%%%%%%%%%%%%%%%%%%%%%%%%


\clearpage
\thispagestyle{empty}

\newcommand*{\myhref}[2]{%
\hypersetup{urlcolor=myred}\href{#1}{#2}}

\makeatletter
\setlength{\@fpbot}{0pt}
\makeatother

\newcolumntype{M}{@{}>{\columncolor{white}[0pt][0pt]}c}

\begin{table}
\def\arraystretch{1.2}
\begin{tabular}{c}
\setlength{\extrarowheight}{7pt}

\begin{minipage}[c]{16.7cm}
    \begin{tabular}{c}
        \hspace{-0.35cm}
        \includegraphics[width=6cm,valign=b,clip]{esp-logo-standard-horizontal.pdf}
        \vspace{0.6cm}
    \end{tabular}
\end{minipage} \\

\begin{minipage}[l]{16.7cm}
    \linespread{1.2}
    \hyphenpenalty=10000
    \hbadness=10000

    \textbf{\large Disclaimer and Copyright Notice}
    \vspace{0.3em}
    
    \footnotesize
    
    Information in this document, including URL references, is subject to change without notice.
    \vspace{0.3em}
    
    ALL THIRD PARTY'S INFORMATION IN THIS DOCUMENT IS PROVIDED AS IS WITH NO WARRANTIES TO ITS AUTHENTICITY AND \\ACCURACY.
    \vspace{0.3em}
    
    NO WARRANTY IS PROVIDED TO THIS DOCUMENT FOR ITS MERCHANTABILITY, NON-INFRINGEMENT, FITNESS FOR ANY PARTICULAR PURPOSE, NOR DOES ANY WARRANTY OTHERWISE ARISING OUT OF ANY PROPOSAL, SPECIFICATION OR SAMPLE.
    \vspace{0.3em}
    
    All liability, including liability for infringement of any proprietary rights, relating to use of information in this document is disclaimed. No licenses express or implied, by estoppel or otherwise, to any intellectual property rights are granted herein.
    \vspace{0.3em}
    
    The Wi-Fi Alliance Member logo is a trademark of the Wi-Fi Alliance. The Bluetooth logo is a registered trademark of Bluetooth SIG.
    \vspace{0.3em}
    
    All trade names, trademarks and registered trademarks mentioned in this document are property of their respective owners, and are hereby acknowledged.
    \vspace{0.3em}
    
    \textbf{Copyright © \the\year{} Espressif Systems (Shanghai) Co., Ltd. All rights reserved.}
    \vspace{0.3em}

    \textbf{\myhref{https://www.espressif.com/}{www.espressif.com}}
\end{minipage}

\end{tabular}
\end{table}


\end{document}
